\section{Related Work}

\subsection{Demand-Responsive Transit and Dial-a-Ride}

Demand-responsive transit (DRT) systems have been studied extensively under the dial-a-ride problem (DARP) framework, in which passengers submit pickup and delivery requests and a fleet of vehicles must be routed to serve them subject to time-window and capacity constraints. \citet{cordeau2007dial} provide a comprehensive survey of DARP models and exact algorithms. \citet{parragh2008survey} extend this to the broader class of pickup-and-delivery problems. \citet{ho2018survey} survey more recent developments, including dynamic and real-time variants relevant to operational DRT deployments.

\citet{dong2020} extend the DARP to include passenger accept/reject decisions based on service utilities, a formulation closely related to our MNL-based choice model; their work motivates the inclusion of an explicit outside option in the choice model (Section~\ref{sec:mnl}).

The many-to-one variant studied here --- where all passengers share a common destination --- arises naturally in first-mile/last-mile transit and commuter shuttle contexts. Meeting-point strategies, in which passengers walk to a shared pickup location rather than being collected door-to-door, have been shown to substantially improve vehicle efficiency \citep{stiglic2015benefits}. High-capacity ride-sharing with dynamic trip-vehicle assignment further demonstrates that pooling passenger requests at shared stops can achieve near-optimal fleet utilization \citep{alonso2017demand}, while minimum-fleet analysis confirms that smart assignment of shared rides substantially reduces the number of vehicles required to serve a given demand \citep{vazifeh2018addressing}. Our work extends this line by treating the \emph{set} of displayed meeting-point offers as a decision variable rather than a fixed system parameter.

Recent work has extended meeting-point strategies in several directions. \citet{czioska2019} develop methods for identifying real-world meeting points in shared DRT systems using geographic and accessibility criteria. \citet{cortenbach2024} formulate the dial-a-ride problem with meeting points as a mixed-integer program, providing exact solutions for small instances and heuristics for larger ones. \citet{stumpe2024} study the design of taxi ridesharing systems with shared pick-up and drop-off locations, finding that meeting-point flexibility substantially reduces vehicle-kilometres travelled. \citet{zheng2019meeting} analyze meeting points in flex-route transit services and show that introducing passenger walking can improve service productivity under appropriate operating conditions. Our work complements these contributions by focusing on the \emph{display} decision: given a set of feasible meeting-point bundles, which subset should be shown to the passenger?

\subsection{Assortment Optimization under MNL}

The problem of selecting which products to display to a customer who chooses according to an MNL model is known as assortment optimization. \citet{talluri2004revenue} establish that revenue-ordered assortments are optimal under the uncapacitated MNL model and derive the Lambert-W pricing formula that we adopt in Section~\ref{sec:mnl}. \citet{rusmevichientong2010dynamic} extend this to capacity-constrained and dynamic settings, showing that greedy revenue-ordered policies remain near-optimal under mild conditions.

\citet{wang2012} extend the assortment problem to the capacitated setting, showing that the revenue-ordered property breaks down when products have heterogeneous costs --- a result directly relevant to our setting where bundles have heterogeneous operational costs. \citet{desir2022capacitated} establish hardness and approximation results for capacitated assortment optimization, further illustrating why constrained assortment problems generally require more than simple ordering rules.

Our setting differs from classical assortment optimization in two important respects. First, the ``products'' (service bundles) are generated dynamically per request and their attributes are predicted rather than known. Second, the operator's objective couples menu design with vehicle routing costs, so the revenue-maximising assortment and the profit-maximising menu can diverge. The menu-value surrogate in Eq.~\eqref{eq:menu_value} is designed to capture this coupling through the system evaluation cost $\hat{c}_{\text{sys}}(b)$.

\subsection{Ride-Sharing and On-Demand Mobility}

\citet{agatz2012optimization} survey optimization approaches for dynamic ride-sharing, highlighting the tension between real-time assignment and service quality. \citet{zha2018surge} analyze how pricing mechanisms affect driver supply and passenger demand in ride-sourcing, showing that platform-level price adjustments have heterogeneous effects on service quality and operator revenue. \citet{fielbaum2021ondemand} study on-demand ridesharing with optimized pickup and drop-off walking locations, a setting closely related to ours; their results confirm that allowing flexible pickup locations significantly expands the feasible service region and improves system efficiency.

Our contribution is complementary: rather than optimizing the \emph{location} of meeting points, we study which subset of pre-generated meeting-point offers to \emph{display} to the passenger, and how the display decision interacts with MNL-based pricing and passenger choice.

\subsection{Vehicle Routing and Route Optimization}

The back-end route optimizer in our system uses the Hybrid Genetic Search (HGS) algorithm of \citet{vidal2022hybrid}, which provides high-quality solutions to capacitated vehicle routing problems with time windows. We treat the routing layer as a black box and focus on the menu-policy layer; the separation between these layers is a deliberate design choice that allows menu policies to be evaluated without re-solving the full VRPTW at every decision epoch.

\subsection{Positioning of This Work}

The present paper sits at the intersection of assortment optimization and DRT operations. Relative to prior work, it brings together the menu-display decision in a many-to-one DRT system with a three-layer framework combining MNL-based passenger choice, common-offset pricing, route-aware menu-value evaluation, and exact-small/greedy-large assortment solving. The paired methodology separates approximation diagnostics, robust filtering behavior, uptake-regime sensitivity, and cross-instance directional checks, making the filtering layer and the greedy approximation empirically auditable rather than hidden implementation details. The tension between platform profit, pooling decisions, pricing, and passenger welfare that motivates our uptake-regime analysis has been noted in ride-sharing equilibrium models \citep{zhang2021pool}, and stated-preference and hybrid-choice approaches to demand-adaptive and on-demand transit service design \citep{frei2017flexing,alsaleh2023ondemand} underscore the importance of modeling passenger choice explicitly when evaluating menu policies.
